In this subsection, which is preparatory to all later material, we give a more comprehensive introduction to the category of Artin--Tate motivic spectra over a field. This begins with introducing the smaller subcategories of Artin and Tate objects, which together generate the category of Artin--Tate objects. The bulk of the subsection is spent building a collection of important examples of Artin--Tate motivic spectra. Over $\R$ a special role is played by the invertible objects and so we focus specific attention there. We close the section by discussing homotopy groups.

\begin{dfn}
  The categories of Artin, Tate and Artin--Tate motivic spectra over $k$ are the stable, full subcategories of $\SHk$, closed under tensor products and colimits, with the following generators:
  \begin{itemize}
  \item The category of Artin motivic spectra, $\SHka$, is generated by the motives of finite \'etale $k$-algebras.
  \item The category of Tate motivic spectra, $\SHkt$, is generated by the motives of $\P_k^1$ and $(\P_k^1)^{-1}$ \footnote{This category is sometimes referred to as the cellular category following \cite{DICell}.}.
  \item The category of Artin--Tate motivic spectra, $\SHkat$, is generated by the motives of finite \'etale $k$-algebras together with $\P_k^1$ and $(\P_k^1)^{-1}$.
  \end{itemize}
  Since these subcategories are closed under tensor products and colimits, they each inherit the structure of a stable, presentably symmetric monoidal category from $\SHk$. \footnote{In the sequel, we will state without remark results that are only known to be true after the characteristic of $k$ is inverted. The reader is invited to restrict themself to $k$ of characteristic zero if they prefer.}
\end{dfn}

We now turn to examples of objects in each of these categories, starting with the trivial and heading towards the non-trival.

\begin{exm} \label{exm:Sp-unit}
  Every stable, presentably symmetric monoidal category admits a unique symmetric monoidal functor from the category of spectra \cite[Corollary 4.8.2.19]{HA} \footnote{Another way of saying this is that $\Sp$ is the initial object of $\mathrm{CAlg}(Pr^{L,\mathrm{stab}})$.}.
  This gives us objects $X \otimes \o_k$ for every spectrum $X$.
  Of particular interest are the integer simplicial suspensions of the unit, $\Ss^{n} \otimes \o_k$.
\end{exm}

Inductively applying the homotopy purity theorem \cite[Theorem 3.2.23]{MorelVoevodsky} to the decomposition $\P^n _k = \mathbb{A}^{n} _k \coprod \P^{n-1} _k$, we learn that:

\begin{exm}[{\cite[Example 2.12]{DICell}}]
  For each $n$, the motivic spectrum associated to $\P^n_k$ is Tate.
\end{exm}

More generally, using a Bialynicki-Birula decomposition, Wendt shows that any smooth projective variety which admits a $\G_m$-action whose fixed points are discrete and rational is Tate \cite{Wendt} \footnote{In the case where the fixed points are not rational, the authors wonder what conditions are necessary to guarantees the motive is Artin--Tate.}.

\begin{exm}[{\cite[Theorem 6.4]{DICell}, \cite[Proposition 8.1]{HM}, \cite[Proposition 8.12]{etaPeriodic}}]
  Each of the following commutative algebras is Tate:
  \[ \MGL,\quad \kq,\quad \kgl,\quad \mathrm{M}\Z,\quad \mathrm{M}\F_p. \]
\end{exm}

\begin{exm} \label{exm:galois-corr}
  Essentially by definition, the Galois correspondence provides functors
  \begin{center}
    \begin{tikzcd}
      {\left\{ \footnotesize\parbox{2.25cm}{finite, continuous Gal($\overline{k}/k$)-sets} \right\}} \ar[r] &
      {\left\{ \footnotesize\parbox{1.5cm}{finite, etale $k$-algebras} \right\}^{op}} \ar[r, hook] \ar[d] &
      \mathrm{Sm}_k \ar[d] \\
      & \SHka \ar[r, hook] &
      \SHk,
    \end{tikzcd}
  \end{center}
  which give our first examples of Artin objects.
\end{exm}

Since the functor in the example above sends disjoint unions to sums there is no loss in generality if we restrict to transitive Gal($\overline{k}/k$)-sets (field extensions). In the case of $\R$ this gives only two objects: $\Spec(\R)$, which we will temporarily denote $\o_\R$ since is it the monoidal unit of the category, and $\Spec(\C)$. Using the natural map we get a cofiber sequence
\begin{align} \Spec(\C) \to \o_\R \xrightarrow{a} \Ss^{\C}. \label{eq:pic-C} \end{align}
On the level of $\C$-points (with Galois action), this cofiber sequence gives the representation sphere $\Ss^{\sigma}$.




Given a quadric in the plane, $V$, we can take its closure in $\P_k^2$ to obtain $\overline{V}$ which is a form of $\P_k^1$. If we let $x_1,\dots,x_n$ denote the points at infinity, then by the homotopy purity theorem \cite[Theorem 3.2.23]{MorelVoevodsky} we have a cofiber sequence
\[ V \to \overline{V} \to \bigoplus_{i=1,\dots,n} \P_{k(x_i)}^1. \]
Under the assumption that $V$ admits a rational point, $\overline{V}$ is just $\P^1_k$, and so we obtain our first non-trivial example of a motive which is neither Artin nor Tate, but is Artin--Tate:

\begin{exm}
  The motive of an affine quadric which admits a rational point is Artin--Tate.
\end{exm}

Over $\R$ there is a particular affine quadric of interest to us.
It is the object $Q$ which appeared in the statement of \Cref{thm:main}.

\begin{exm}
  Let $Q \coloneqq \{x^2+y^2=1\} \subset \AA^2 _{\R}$, which we shall call the algebraic circle.
  %% From the previous example we know that $Q$ is Artin--Tate.
  $Q$ has the additional property that it is a form of $\G_m$ \footnote{By Cartier duality, forms of $\G_m$ are classified by rank 1 lattices with Galois action. $Q$ corresponds to the unique nontrival action.}.
  Its group scheme structure comes from the usual rule for multiplication of complex numbers,
  namely
  \[ (x_1,y_1) \cdot (x_2,y_2) = (x_1 x_2 - y_1 y_2,\ x_1 y_2 + y_1 x_2). \]
  %% Furthermore, it is a form of $\G_m$---there is an isomorphism $(Q)_\C \cong (\G_m)_\C$ of commutative group schemes over $\C$.

  In \Cref{sec:ta} we will construct the maps $\ta: \Ss^1 \to Q_p$ using our explicit understanding of this group scheme.
\end{exm}

\begin{prop}[{\cite[Proposition 1.1]{HuPicard}}] \label{prop:pic-Q}
  There is an equivalence $Q \otimes \Ss^\C \simeq \P^1 _\R$.
  In particular, both $\Ss^\C$ and $Q$ are invertible.
\end{prop}

At this point we are ready to define the appropriate notion of homotopy groups for studying $\SHR^{\at}$.
In the case of $\SHC^{\at}$, the category has a bi-graded family of compact invertible generators given by the spheres $\Ss^{p - 2w} \otimes (\P_\C^1)^{\otimes w}$. Therefore, it is typical to study objects at the level of their bigraded homotopy groups, given by
\[ \pi_{p,w}^\C(X) \coloneqq \pi_0\Map ( \Ss^{p-2w} \otimes (\P_\C^1)^{\otimes w}, X ). \]
Similarly, $\SHRt$ has a bi-graded family of compact invertible generators given by the spheres $\Ss^{p-2w} \otimes (\P_\R^1)^{\otimes w}$, and it is typical to study objects through their bi-graded homotopy groups.
Using \Cref{prop:pic-Q} and \Cref{eq:pic-C}, the category $\SHR^{\at}$ has a tri-graded family of compact invertible generators given by the spheres $\Ss^{p-w} \otimes (\Ss^\C)^{\otimes q-w} \otimes (\P_\R^1)^{\otimes w}$, and we will study most objects through their tri-graded homotopy groups.

\begin{ntn}
  We endow the categories $\Sp$, $\Sp_{C_2}$, $\SHC$ and $\SHR$ with Picard gradings by spheres as follows:
  \begin{center}
    \renewcommand{\arraystretch}{1.1}
    \begin{tabular}{|l|l|l|} \hline
      Category & Picard & Spheres \\\hline\hline
      $\Sp$ & $\Z$ & $\Ss^p \simeq (S^1)^{\otimes p}$ \\\hline
      $\Sp_{C_2}$ & $\Z \times \Z$ & $\Ss^{p+q\sigma} \simeq (S^1)^{\otimes p} \otimes (S^\sigma)^{\otimes q}$ \\\hline
      $\SHC$ & $\Z \times \Z$ & $\Ss^{p,w} \simeq (S^1)^{\otimes p-2w} \otimes (\P_\C^1)^{\otimes w}$ \\\hline
      $\SHR$ & $\Z \times \Z \times \Z$ & $\Ss^{p,q,w} \simeq (S^1)^{\otimes p-w} \otimes (S^\C)^{\otimes q-w} \otimes (\P_\R^1)^{\otimes w}$ \\\hline
    \end{tabular}
    \renewcommand{\arraystretch}{1.0}
  \end{center}
  
  For any $p,q,w \in \Z$, and any $\R$-motivic spectrum $X$, we let $\pi^{\R}_{p,q,w} X$ denote the group of homotopy classes of maps,
  \[ (S^1)^{\otimes p-w} \otimes (S^\C)^{\otimes q-w} \otimes (\P_\R^1)^{\otimes w} \to X. \]
\end{ntn}

\begin{exm}
  For convenience we record what the invertible objects considered in this section look like under this new notation:
  \begin{center}
    \begin{tabular}{ccc}   
      $\Ss^{0,0,0} \cong \o_\R$ \qquad\qquad\qquad &
      $\Ss^{1,0,0} \cong S^1$ \qquad\qquad\qquad & 
      $\Ss^{0,1,0} \cong S^{\C}$ \\
      $\Ss^{0,1,1} \cong \G_m$ \qquad\qquad\qquad &
      $\Ss^{1,0,1} \cong Q$ \qquad\qquad\qquad &
      $\Ss^{1,1,1} \cong \P_\R^1$
    \end{tabular}
  \end{center}
  In the case of $\Ss^{0,0,0}$ we will often drop the indices for brevity, writing only $\Ss$.
  The two maps between Picard elements considered thus far, $\ta$ and $a$, become
  \[ \ta \in \pi_{0,0,-1}^\R \Ss_p \qquad\quad \text{ and } \qquad\quad a \in \pi_{0,-1,0}^\R \Ss. \]
  %% In the case when $X$ is the unit we will often omit it from the notation for brevity.
\end{exm}

\begin{rmk}
  The tri-graded spheres in $\SHR^{\at}$ which are Tate are precisely the spheres of the form $\Ss^{p,q,q}$.
  For this reason the Tate category only sees a `slice' of the total information present in the tri-graded homotopy groups.

  The tri-graded spheres in $\SHR^{\at}$ which are Artin are precisely the spheres of the form $\Ss^{p,q,0}$.
  The Artin category similarly only sees a `slice' of the total homotopy groups.

  %% Traditionally, $\SHR$ is viewed as bigraded by the spheres $\Ss^{p,w} \simeq (S^1)^{\otimes p-w} \otimes (\G_m)^{w}$. The comparison with our trigraded spheres is given by the equivalence $\Ss^{p,w} \simeq \Ss^{p-w,w,w}$, and from this we see that the bigraded sheres capture precisely those trigraded spheres $\Ss^{p,q,w}$ satisfying $q=w$.
\end{rmk}

%% as a  defined by [people]. We regard this category as the correct analog of the category over cellular objects over an algebraically closed field and propose that it may be possible to extend the full package of results known in the cellular case over algebraically closed fields to the Artin-Tate category over general fields. 


%% the the Picard elements $\Ss^1$, $\G_m$, $\Ss^{\C}$ and their inverses.

%% Let $\Ss_2^{\wedge}$ denote the $2$-completed unit in $\SHR$, and consider the category $\Ss_2^{\wedge}$-mod of $\Ss_2^{\wedge}$-module spectra.  We similarly define $\SHgc_2$ to be the smallest full subcategory of $\Ss_2^{\wedge}$-mod, closed under colimits and tensor products over $\Ss^{\wedge}_{2}$, containing the Picard elements $(\Ss^1)^{\wedge}_2, (\G_m)_2^{\wedge}$, $(\Ss^{\C})_2^{\wedge}$, and their inverses.

%% Both $\SHgc$ and $\SHgc_2$ are symmetric monoidal stable $\infty$-categories, and by CITE the $2$-completions of these categories are equivalent.  Note however that not every object in $\SHgc_2$ is $2$-complete.
%% \end{dfn}


%% \begin{cnstr}
%% There is a canonical class $\ta: \Ss^{\C} \to (\G_m)^{\wedge}_2$.  We will build $\ta$ in \Cref{sec:tau} via the explicit algebraic geometry of the circle $Q=\Spec \R[x,y]/(x^2+y^2-1)$.  A key point is that $\Sigma^{\infty} Q$ is itself a Picard element in $\SHgc$, satisfying a relation $\Ss^{\C} \otimes \Sigma^{\infty} Q \simeq \Ss^1 \otimes \G_m$.
%% \end{cnstr}


%% In this section, we introduce some exotic Picard elements in $\SHR$, use them to define trigraded homotopy groups, and define some maps between them, in particular the element $\ta \in \pi^{\R} _{0,0,-1} \Ss_p$, which serves as a Picard-graded analogue of the class $\tau \in \pi^{\C} _{0,-1} \Ss_p$.

%% \todo{Give a general definition of things}
%% \todo{Conjecture: every pic element is already in this category}

%% We discuss Picard gradings in $\Sp$, $\Sp_{C_2}$, $\SHC$ and $\SHR$. We begin by introducing the invertible objects that will generate our Picard gradings.

%% \begin{rec}
%%   Given an category $\CC$ admitting colimits, we let $S^1$ denote the simplicial circle, i.e. the colimit of the constant diagram on $S^1$ at the terminal object.
%%   In $\mathrm{Spaces}_{C_2}$, let $S^{\sigma}$ denote the unreduced suspension of $C_2$. Similarly, in $\MR$, we let $S^{\C}$ denote the unreduced suspension of $\Spec \C$.
%%   Finally, in $\MR$ and $\MC$, we let $\G_m$ denote the multiplicative group scheme.
%%   By abuse of notation, we will use the same notation for the suspension spectra of $S^1$, $S^{\sigma}$, $S^\C$ and $\G_m$.

%%   Then $S^1$ is invertible in any stable category.
%%   It is well-known that $S^{\sigma}$ is invertible in $\Sp_{C_2}$, and that $\G_m$ is invertible in $\SHC$ and $\SHR$.
%%   It is a result of Hu that $S^\C$ is invertible in $\SHR$ \cite[Proposition 1.1]{HuPicard}.
%% \end{rec}


%
%\begin{rec}
  %Let $\sigma$ denote the real sign representation of $C_2$, and let $S^\sigma$ denote the one-point compactification of $\sigma$. There is also a cofiber sequence of pointed $C_2$-spaces $(C_2)_+ \to S^0 \to S^{\sigma}$.
%
  %The stabilized sphere $\Ss^{\sigma}$ is invertible, and one defines $\Ss^{p+q\sigma} = (\Ss^1)^{\otimes p} \otimes (\Ss^\sigma)^{\otimes q}$. This endows the category of $C_2$-spectra with an $RO(C_2)$-grading, which may also be viewed as a bigrading.
%\end{rec}
%
%\begin{rec}
  %Let $S^1 \in \MC$ denote the simplicial circle and let $\G_m \in \MC$ be the motivic space associated to the multiplicative group. Let $\Ss^{1,0} \coloneqq \Sigma^\infty S^1$ and $\Ss^{1,1} \coloneqq \Sigma^\infty \G_m$. We then define $\Ss^{a,b} = (\Ss^{1,0})^{\otimes a-b} \otimes (\Ss^{1,1})^{\otimes b}$. This endows $\SHC$ with a bigrading.
%\end{rec}


%% \begin{rmk}
%%   Traditionally, $\SHR$ is viewed as bigraded by the spheres $\Ss^{p,w} \simeq (S^1)^{\otimes p-w} \otimes (\G_m)^{w}$. The comparison with our trigraded spheres is given by the equivalence $\Ss^{p,w} \simeq \Ss^{p-w,w,w}$, and from this we see that the bigraded sheres capture precisely those trigraded spheres $\Ss^{p,q,w}$ satisfying $q=w$.
%% \end{rmk}

%% There is one more explicit Picard element in $\SHR$ that shall be of interested to us.

%% \begin{prop}
%%   Stably, there is an equivalence $Q \simeq \Ss^{1,0,1}$. In particular, $Q$ is invertible.
%% \end{prop}

%% \begin{proof}
%%   Hu has shown that $Q \otimes \Ss^\C \simeq \P^1 _\R$ \cite[Proposition 1.1]{HuPicard}. On the other hand, it is standard that $\G_m \otimes \Ss^1 \simeq \P^1 _\R$, from which it follows that $\P^1 _\R \simeq \Ss^{1,1,1}$, which implies the desired result.
%% \end{proof}





%Traditionally, $\SHR$ is also studied via its bigraded homotopy groups. In this paper, we will endow $\SHR$ further with a trigrading.
%As our first step in introducing this trigrading, we recall several invertible $\R$-motivic spectra.
%
%\begin{dfn}
  %We define several $\R$-motivic spaces as follows:
  %\begin{itemize}
    %\item $S^1$ is the simplicial circle.
    %\item $\G_m$ is the multiplicative group.
    %\item $\P^1$ is the projective line.
    %\item $S^\C$ is the unreduced suspension of $\Spec \C$.
    %\item $Q$ is the affine quadric $\{x^2 + y^2 = 1\} \subseteq \AA^2$, which is an abelian group scheme under complex multiplication: $(x_1,y_1) \cdot (x_2,y_2) = (x_1 x_2 - y_1 y_2, x_1 y_2 + y_1 x_2)$.
  %\end{itemize}
  %We will abuse notation by using the same notation to denote the associated $\R$-motivic suspension spectra.
%\end{dfn}
%
%\begin{rmk}
  %Note that $Q$ is a form of $\G_m$: there is an isomorphism of group schemes $Q_\C \cong (\G_m)_\C$. Moreover, there is an equivalence of $\C$-motivic spaces $(S^\CC)_\CC \simeq (S^1)_\CC$.
%\end{rmk}
%
%\begin{prop}
  %The $\R$-motivic spectra $S^1$,$\G_m$, $\P^1$, $S^\C$ and $Q$ are all invertible. Moreover, there are relations $S^1 \otimes \G_m \simeq S^\C \otimes Q \simeq \P^1$.
%\end{prop}
%
%\begin{proof}
  %It is classical that $S^1$, $\G_m$ and $\P^1$ are invertible, as is the relation $S^1 \otimes \G_m \simeq \P^1$. The relation $S^\C \otimes Q \simeq \P^1$ is \cite[Proposition 1.1]{HuPicard} and implies that $S^\C$ and $Q$ are invertible.
%\end{proof}
%
%The $\R$-motivic spectra $S^1, \G_m$ and $\P^1$ are inverible by the definition of $\SHR$, and there is the classical relation $S^1 \otimes \G_m \simeq \P^1$ among them. In \cite{HuPicard}, it is shown that $S^\C \otimes Q \simeq \P^1$, showing that $S^\C$ and $Q$ are also invertible. After base change to $\C$, there are equivalences $S^{\C} _{\C} \simeq S^1 _\C$ and $Q_{\C} \cong (\G_m)_\C$, the latter originating from an isomorphism of group schemes.
%
%Using the above Picard elements and relations, we are able to introduce a trigrading on the $\R$-motivic stable homotopy category, extending the usual bigrading.
%
%\begin{ntn}
  %We define trigraded spheres as follows. First, we define:
  %\begin{itemize}
    %\item $\Ss^{1,0,0} \coloneqq S^1$.
    %\item $\Ss^{0,1,0} \coloneqq S^{\C}$.
    %\item $\Ss^{0,1,1} \coloneqq \G_m$.
    %\item $\Ss^{1,0,1} \coloneqq Q$.
    %\item $\Ss^{1,1,1} \coloneqq \P^1$.
  %\end{itemize}
  %In general, we define $\Ss^{p,q,w}$ for $(p,q,w) \in \Z^{\times 3}$ by demanding that $\Ss^{p_1,q_1,w_1} \otimes \Ss^{p_2,q_2,w_2} \cong \Ss^{p_1+p_2,q_1+q_2,w_1+w_2}$. This is possible because of the equivalences $S^1 \otimes \G_m \simeq S^\C \otimes Q \simeq \P^1$.
%\end{ntn}
%
%\begin{rmk}
  %Recall the usual motivic bigrading, where $\Ss^{1,0} \coloneqq S^1$ and $\Ss^{1,1} \coloneqq \G_m$. Then the bigraded sphere $\Ss^{a,b}$ is equivalent to the trigraded sphere $\Ss^{a-b,b,b}$. A trigraded sphere $\Ss^{p,q,w}$ is equivalent to a bigraded sphere precisely when $q=w$.
%\end{rmk}
%
%Under this trigrading, the first two gradings keep track of the underlying $RO(C_2)$-grading, and the third grading is the motivic weight.
%
%\begin{prop}
  %Under the Betti realization functor $\b : \SHR \to \Sp_{C_2}$, we have $\b(\Ss^{a,b,c}) \simeq \Ss^{a+b\sigma}$.
  %Under the base change functor $(-)_\C : \SHR \to \SHC$, we have $(\Ss^{a,b,c})_\C \cong \Ss^{a+b,c}$. \todo{What's a better way to denote this?}
%
  %As a consequence, we find that $\Ss^{p_1,q_1,w_1} \simeq \Ss^{p_2,q_2,w_2}$ implies that $(p_1,q_1,w_1) = (p_2,q_2,w_2)$.
%\end{prop}

%% The categories $\Sp$ and $\Sp_{C_2}$ are generated under colimits by the Picard spheres generated above. This is not true of $\SHC$ and $\SHR$, so if we want to study these categories via the Picard graded homotopy groups indexed by the spheres considered above, we would do well to restrict ourselves to the full subcategories generated by these spheres.

